\section{Budget-Constrained Proof Search}

\subsection{Interaction Model}

Let \(x\) be a theorem-proving task and \(s_t\) the controller state after
\(t\) computations. The state contains a verifier-grounded proof workspace,
the action history, and a remaining resource vector. At each step the
controller selects an action \(a_t \in \mathcal{A}(s_t)\), observes an outcome
\(o_t\), and applies a transition
\[
    s_{t+1} = T(s_t,a_t,o_t).
\]
The transition is deterministic given the action and tool result; uncertainty
enters through model generations, retrieval outcomes, and the unknown
repairability of a branch.

Our action space includes proof implementation, local repair, retrieval,
capability testing, decomposition, argument refinement, representation change,
model escalation, branch pruning, assembly, and stopping. An action has a
nonnegative resource cost vector
\[
  \mathbf{c}(s,a) =
  (c_{\mathrm{in}}, c_{\mathrm{out}}, c_{\mathrm{model}},
   c_{\mathrm{check}}, c_{\mathrm{retrieve}}, c_{\mathrm{time}},
   c_{\mathrm{money}}).
\]
The run is feasible when cumulative cost is component-wise bounded by a budget
\(\mathbf{B}\). We retain the vector for reporting and enforce hard limits
directly. For ranking actions, a deployment-specific nonnegative weight vector
\(\boldsymbol{\lambda}\) converts incremental cost to
\[
  \bar c(s,a)=\boldsymbol{\lambda}^{\top}\mathbf{c}(s,a).
\]
This separation prevents a cheap monetary price from concealing excessive
checker calls or latency.

Let \(Y_\pi(x,\mathbf{B})\in\{0,1\}\) indicate whether policy \(\pi\) produces
a proof accepted by Lean and by the deterministic safety review before budget
exhaustion. The primary objective is
\[
  \max_\pi\ \mathbb{E}_{x\sim\mathcal{D}}
  [Y_\pi(x,\mathbf{B})],
\]
evaluated over a range of budgets rather than at one arbitrary cutoff. A
secondary objective minimizes expected cost conditional on success.

\subsection{Local Cost-Sensitive Decisions}

The simplest case already yields a useful policy target. Suppose candidate
action \(i\) costs \(c_i>0\), succeeds with probability \(p_i\), outcomes are
conditionally independent, and checking stops after the first success.

\begin{theorem}[Candidate ordering]
Among fixed candidate orderings, sorting by nonincreasing \(p_i/c_i\)
minimizes expected cost to the first success, with zero-probability candidates
last.
\end{theorem}

\begin{proof}
Consider adjacent candidates \(i\) and \(j\) after a common prefix. Trying
\(i\) first has conditional expected cost
\(c_i+(1-p_i)c_j\); trying \(j\) first costs
\(c_j+(1-p_j)c_i\). The first ordering is no worse exactly when
\(p_i/c_i \ge p_j/c_j\). Repeated adjacent swaps remove all inversions.
\end{proof}

This result does not assume that proof success probabilities are easy to
estimate. It identifies the quantity that a controller should approximate from
verifier feedback, action type, branch history, and model tier.

Now consider a cheap action of cost \(c_L\) that solves the task with
probability \(p_L\), followed on failure by a strong-model action of cost
\(C_H\).

\begin{proposition}[Cheap search before escalation]
Cheap-first execution has lower expected cost than direct escalation whenever
\[
    c_L < p_L C_H.
\]
\end{proposition}

Cheap search can also improve the state even when it does not finish the
proof. If direct strong generation succeeds with probability \(p_H(s_0)\),
whereas cheap search moves failures to \(s_1\), the cheap-first policy has
lower expected cost per solved task when
\[
\frac{c_L+(1-p_L)C_H}
     {p_L+(1-p_L)p_H(s_1)}
<
\frac{C_H}{p_H(s_0)}.
\]
This condition captures two distinct benefits: direct cheap success and
verifier-grounded state improvement before escalation.

\subsection{Residual Verifier Progress}

Absolute LLM value judgments are often poorly calibrated. We instead derive a
progress feature vector \(\phi(s)\) from the verifier and workspace. Components
include accepted obligations, newly ready dependents, diagnostic transitions,
goal fingerprints, stalled streaks, and verified-prefix information. The
residual progress of an action is
\[
  \Delta(s,a,s') = \mathbf{w}^{\top}
  \bigl(\phi(s')-\phi(s)\bigr).
\]
When \(\phi\) is only the terminal acceptance indicator,
\(\mathbb{E}[\Delta]/c\) reduces exactly to \(p/c\). Richer features provide a
denser but still verifier-grounded signal. We treat the relationship between
these proxy features and eventual success as an empirical question rather
than assuming that fewer goals always means an easier proof.
